\section{Related Work}
\paragraph{Persistent Memory and RAG.}
Retrieval-augmented generation improves factuality and coverage by accessing external knowledge at inference time~\citep{lewis2020retrieval}. Flat nearest-neighbor memory has likewise been used to improve language-model behavior through direct memory lookup~\citep{khandelwal2020generalization}. In agent settings, memory is often maintained as a chronological stream of observations, reflections, or summaries~\citep{park2023generative}. These approaches demonstrate the value of persistent memory but treat it as a flat collection, without exploiting the event and topic structure inherent in conversation.

\paragraph{RL for Search and Memory.}
Search-R1 treats external search as a policy decision and optimizes reasoning-search trajectories with relative reinforcement learning~\citep{jin2025search}. Memory-R1 extends this perspective to persistent memory, training a Memory Manager and an Answer Agent using downstream question-answering reward~\citep{yan2025memory}. These works establish the key principle that information access and memory manipulation should be trained from task outcomes rather than from heuristic rules. However, Search-R1 addresses external search rather than persistent conversational memory, and Memory-R1 operates over flat memory without addressing tree-structured access or updates. Our work inherits this outcome-driven perspective while extending it to tree-structured memory with learned access and update policies.

\paragraph{Structured Memory Access.}
Structured memory representations expose relations, hierarchies, and subtrees that are unavailable in flat banks. Semantic XPath is representative of work that leverages this structure at inference time by generating path-like queries over memory trees~\citep{semanticxpath2026}. This line of work shows that hierarchical organization can improve memory access, but it typically assumes the structure already exists and relies on heuristic or prompted query generation. Our setting differs in two respects: the tree is induced incrementally from raw conversation, and access over it is optimized as a learned policy using downstream reward.

\paragraph{Gap addressed by this work.}
Prior work treats memory structure, learned access, and learned updates as separate concerns. \textsc{Struct Memory-R1} addresses all three within a unified framework.
